\documentclass[12pt]{article}
\usepackage{amsmath, amssymb, amsthm, bm}
\usepackage[margin=1in]{geometry}
\usepackage[skip=0.6em]{parskip}

\newtheorem{theorem}{Theorem}
\newtheorem{proposition}{Proposition}
\newtheorem{lemma}{Lemma}
\newtheorem{corollary}{Corollary}
\theoremstyle{remark}
\newtheorem{remark}{Remark}

\begin{document}

\begin{center}
{\Large Feasible normals and admissibility\\ in the restricted bivariate normal problem\\ under general correlation}\\[3pt]
{\small Ivan A. Canay :: August 27, 2026}
\end{center}

\medskip

\section*{Setup and statement}

Let
\[
X\sim N_2(\mu,\Omega_\rho),
\qquad
\Omega_\rho=
\begin{pmatrix}
1&\rho\\
\rho&1
\end{pmatrix},
\qquad -1<\rho<1~,
\]
and define
\[
\Lambda=\{\mu:\mu_2\le \mu_1\}~,
\]
\[
\Lambda_0=\{\mu\in\Lambda:\mu_1\le0\}~,
\qquad
\Lambda_1=\{\mu\in\Lambda:\mu_1>0\}~.
\]
A nonrandomized test is written as
\[
\phi_A(x)=I\{x\notin A\}~,
\]
where $A\subseteq\mathbf R^2$ is the acceptance region.  Throughout this note,
$A$ is assumed to be closed, convex, and coordinatewise lower:
\[
x\in A,
\quad y\le x
\quad\Longrightarrow\quad
y\in A~.
\]
For a closed convex set $C$, write
\[
h_C(a):=\sup_{x\in C}a'x
\]
for its support function, allowing the value $+\infty$.

Relative to $P_0=N_2(0,\Omega_\rho)$, the natural parameter corresponding to
the mean $\mu$ is $\vartheta=\Omega_\rho^{-1}\mu$.  A supporting normal $a$
therefore generates the mean ray $\mu_\lambda=\lambda\Omega_\rho a$.  This
suggests the feasible normal cone
\begin{equation}\label{eq:feasible-cone-general}
\begin{split}
\mathcal F_\rho
&:=\{a\in\mathbf R^2:a\ge0,\ \Omega_\rho a\in\Lambda\}\\
&=\{a\in\mathbf R^2:0\le a_2\le a_1\}
=\operatorname{cone}\{(1,0),(1,1)\}~.
\end{split}
\end{equation}
The second equality follows because
\[
(\Omega_\rho a)_1-(\Omega_\rho a)_2
=(1-\rho)(a_1-a_2)~.
\]
Thus the feasible normal cone does not depend on $\rho$, even though the
likelihood-ratio geometry does.  The phrase ``all finite normals lie in
$\mathcal F_\rho$'' means
\begin{equation}\label{eq:finite-normal-condition}
a\ne0,
\quad h_A(a)<\infty
\quad\Longrightarrow\quad
a\in\mathcal F_\rho~.
\end{equation}
For a lower set, finite support already forces $a\ge0$, so failure of
\eqref{eq:finite-normal-condition} is equivalent to the existence of some
$a\ne0$ with
\[
h_A(a)<\infty~,
\qquad
0\le a_1<a_2~.
\]

The main result to be established is the following characterization.

\begin{theorem}[Feasible-normal characterization]\label{thm:iff}
Let $A\subseteq\mathbf R^2$ be closed, convex, and coordinatewise lower, and suppose
\[
P_0(X\in A)=1-\alpha~,
\qquad 0<\alpha<1~.
\]
For every $-1<\rho<1$, the test $\phi_A=I\{X\notin A\}$ is
$\alpha$-admissible against all level-$\alpha$ tests if and only if all finite
normals of $A$ lie in $\mathcal F_\rho$, that is, if and only if
\eqref{eq:finite-normal-condition} holds.
\end{theorem}

The proof has two parts.  The sufficient direction is the Lehmann--Romano ray
argument.  The necessary direction is constructive.  Under nonnegative
correlation, a bad finite normal can be repaired by the same hull and diagonal
calibration used in the identity-covariance problem.  Under negative
correlation, the proof first checks whether $A$ is lower along the second
extreme ray of the transfer order.  If it is not, a conditional monotone
rearrangement strictly improves the test.  If it is, the resulting slope bound
restores the gain-loss ordering needed for the diagonal repair.

\section*{1. Sufficiency}

We use the following standard consequence of the Lehmann--Romano complete-class
result for exponential families.

\begin{proposition}[Lehmann--Romano ray criterion]\label{prop:LR}
Consider an exponential family with natural statistic $T$ and natural parameter
space $\Theta$.  Let $\Theta_H$ be the null parameter set and let $\Theta'$ be
an alternative set.  Let $A_0$ be a closed convex acceptance region for a
nonrandomized test.  Suppose that whenever
\[
A_0\subseteq\{t:a't\le c\}
\]
for some nonzero $a$, there exist $\vartheta^*\in\Theta$ and
$\lambda_n\to\infty$ such that
\[
\vartheta^*+\lambda_n a\in\Theta'~.
\]
Then the test with acceptance region $A_0$ is $d$-admissible.  If, in addition,
the test has exact size $\alpha$ attained at a finite null point, then it is
$\alpha$-admissible.  These conclusions allow randomized competing tests.
\end{proposition}

\begin{lemma}[Sufficiency of feasible normals]\label{lem:sufficiency}
If $A$ is closed, convex, coordinatewise lower, satisfies
$P_0(A)=1-\alpha$, and satisfies \eqref{eq:finite-normal-condition}, then
$\phi_A$ is $\alpha$-admissible against all level-$\alpha$ tests.
\end{lemma}

\begin{proof}
Relative to $P_0$, the normal experiment has the canonical representation
\[
\frac{dP_\mu}{dP_0}(x)
=\exp\left\{(\Omega_\rho^{-1}\mu)'x
-\frac12\mu'\Omega_\rho^{-1}\mu\right\}~,
\]
with natural parameter $\vartheta=\Omega_\rho^{-1}\mu$ and statistic
$T(x)=x$.  The alternative set in the natural parameter space is
\[
\Theta'=\Omega_\rho^{-1}\Lambda_1~.
\]

We verify the ray condition in Proposition~\ref{prop:LR}.  Suppose
\[
A\subseteq\{x:a'x\le c\}
\]
for some nonzero $a$.  Then $h_A(a)<\infty$, so by
\eqref{eq:finite-normal-condition}, $a\in\mathcal F_\rho$.  Consider the
natural-parameter ray $\vartheta_\lambda=\lambda a$.  Its corresponding mean is
\[
\mu_\lambda=\lambda\Omega_\rho a~.
\]
By \eqref{eq:feasible-cone-general}, $\mu_\lambda\in\Lambda$.  Moreover,
$a\ne0$ and $0\le a_2\le a_1$ imply $a_1>0$, and
\[
(\mu_\lambda)_1
=\lambda(a_1+\rho a_2)>0~.
\]
Indeed, this is immediate when $\rho\ge0$, while for $\rho<0$,
\[
a_1+\rho a_2\ge (1+\rho)a_1>0~.
\]
Hence $\mu_\lambda\in\Lambda_1$, or equivalently
$\vartheta_\lambda\in\Theta'$, for every $\lambda>0$.  The ray condition holds
with $\vartheta^*=0$ and any sequence $\lambda_n\to\infty$.

Because $A$ is coordinatewise lower, the rejection region $A^c$ is
coordinatewise upper.  For every $\mu\in\Lambda_0$, we have $\mu\le0$.
Coupling $X=\mu+Z$ with $Z\sim N_2(0,\Omega_\rho)$ gives
$\mu+Z\le Z$ coordinatewise, and therefore
\[
P_\mu(X\notin A)\le P_0(Z\notin A)=\alpha~.
\]
Thus the test has level $\alpha$, and exact size is attained at the finite null
point $\mu=0$.  Proposition~\ref{prop:LR} gives $\alpha$-admissibility against
all level-$\alpha$ tests, including randomized competitors.
\end{proof}

\section*{2. The transfer order}

The necessity proof uses the order with respect to which every
restricted-alternative likelihood ratio is increasing.  Define
\[
T_{1,\rho}(x):=(\Omega_\rho^{-1}\mathbf 1)'x
=\frac{x_1+x_2}{1+\rho}~,
\qquad
T_{2,\rho}(x):=-(\Omega_\rho^{-1}e_2)'x
=\frac{\rho x_1-x_2}{1-\rho^2}~,
\]
and write
\[
x\succeq_{\Lambda,\rho}y
\quad\Longleftrightarrow\quad
T_{1,\rho}(x)\ge T_{1,\rho}(y)
\quad\text{and}\quad
T_{2,\rho}(x)\ge T_{2,\rho}(y)~.
\]
The positive cone of this order is
\begin{equation}\label{eq:transfer-cone}
\begin{split}
\Pi_\rho
&:=\{d:x+d\succeq_{\Lambda,\rho}x\}\\
&=\operatorname{cone}\{(1,\rho),(1,-1)\}\\
&=\{d:d_1+d_2\ge0,\ \rho d_1-d_2\ge0\}~.
\end{split}
\end{equation}
Intuitively, $\Pi_\rho$ is the cone of universally power-improving moves in
sample space for the restricted alternatives.  The direction $(1,\rho)$ raises
the common-shift score $T_{1,\rho}$ while leaving $T_{2,\rho}$ unchanged.  The
direction $(1,-1)$ raises the slack score $T_{2,\rho}$ while leaving
$T_{1,\rho}$ unchanged.

For $\mu\in\Lambda_1$, write
\[
u:=\mu_1>0~,
\qquad
v:=\mu_1-\mu_2\ge0,
\qquad
\mu=u\mathbf 1-ve_2~.
\]
Then, relative to $P_0$,
\begin{equation}\label{eq:LR-density}
L_\mu(x):=\frac{dP_\mu}{dP_0}(x)
=\exp\left\{uT_{1,\rho}(x)+vT_{2,\rho}(x)
-\frac12\mu'\Omega_\rho^{-1}\mu\right\}~.
\end{equation}
Thus $L_\mu$ is strictly increasing in $T_{1,\rho}$ and weakly increasing in
$T_{2,\rho}$.

We first replace $A$ by its closed $-\Pi_\rho$-hull,
\[
\widehat A_\rho:=\operatorname{cl}(A-\Pi_\rho)~.
\]
Equivalently, $\widehat A_\rho$ is the smallest closed set containing $A$ that
is closed under moves opposite to the restricted likelihood-ratio order.  This
operation removes boundary pieces whose supporting normals point in infeasible
directions.  It may add too much acceptance probability, so the next step will
restore the original size by cutting $\widehat A_\rho$ with a feasible diagonal
halfspace.

For a rectangle, the hull operation retains the diagonal repair from the
identity-covariance problem because $(1,-1)$ remains an extreme ray of
$\Pi_\rho$.  The other extreme ray changes from $(1,0)$ to $(1,\rho)$.  When
$\rho\ge0$, its negative is already a coordinatewise downward move and adds
nothing to a lower acceptance region.  When $\rho<0$, it contributes an
additional likelihood-ratio-improving transfer.  Lemma~\ref{lem:hull} treats
both cases at once.

\begin{lemma}[The feasible-normal hull]\label{lem:hull}
The set $\widehat A_\rho$ is closed, convex, coordinatewise lower, and
$\Pi_\rho$-lower, meaning
\[
x\in\widehat A_\rho~,
\quad p\in\Pi_\rho
\quad\Longrightarrow\quad
x-p\in\widehat A_\rho~.
\]
Moreover, every nonzero finite-support direction of $\widehat A_\rho$ lies in
$\mathcal F_\rho$.  If $A$ has a finite-support direction $a$ with
$a_2>a_1$, then
\[
A\subsetneq\widehat A_\rho
\qquad\text{and}\qquad
P_0(\widehat A_\rho)>P_0(A)~.
\]
\end{lemma}

\begin{proof}
Because $\widehat A_\rho$ is the closure of the Minkowski sum
$A+(-\Pi_\rho)$, it is closed and convex.  It contains $A$ because
$0\in\Pi_\rho$.  It is coordinatewise lower because $A$ is coordinatewise
lower: if $x=y-p$ with $y\in A$ and $p\in\Pi_\rho$, and if $z\le x$, then
\[
z=(y-(x-z))-p~,
\]
where $y-(x-z)\in A$.  Passing to the closure preserves the lower-set
property.  The $\Pi_\rho$-lower property follows from the definition.

The support function satisfies
\[
h_{\widehat A_\rho}(a)=h_A(a)+h_{-\Pi_\rho}(a)~,
\]
with the usual convention for extended-real sums.  Since $-\Pi_\rho$ is
generated by $(-1,-\rho)$ and $(-1,1)$,
\begin{equation}\label{eq:polar-condition}
\begin{split}
h_{-\Pi_\rho}(a)<\infty
&\quad\Longleftrightarrow\quad
a_1+\rho a_2\ge0
\quad\text{and}\quad
a_2\le a_1~.
\end{split}
\end{equation}
Also, since $A$ is coordinatewise lower and nonempty, $h_A(a)<\infty$ implies
$a\ge0$.  If $h_{\widehat A_\rho}(a)<\infty$ and $a\ne0$, it follows that
$a\ge0$ and $a_2\le a_1$.  The remaining inequality in
\eqref{eq:polar-condition} is then automatic: it is immediate for $\rho\ge0$,
while for $\rho<0$,
\[
a_1+\rho a_2\ge(1+\rho)a_1\ge0~.
\]
Therefore
\[
0\le a_2\le a_1~,
\]
so $a\in\mathcal F_\rho$.

Finally, suppose that $h_A(a)<\infty$ for some $a$ with $a_2>a_1$.  Then
$h_{\widehat A_\rho}(a)=+\infty$, because $a$ is not in the polar cone of
$-\Pi_\rho$.  Indeed, $(-1,1)\in-\Pi_\rho$ and
\[
a'(-1,1)=a_2-a_1>0~.
\]
Every point $x\in A$ therefore generates the ray
$x+t(-1,1)\in A-\Pi_\rho$.  Along this ray the support functional tends to
$+\infty$:
\[
a'\{x+t(-1,1)\}=a'x+t(a_2-a_1)\longrightarrow+\infty~.
\]
Hence $\widehat A_\rho$ is not contained in the supporting halfspace
$\{x:a'x\le h_A(a)\}$ that contains $A$, so
$\widehat A_\rho\setminus A$ is nonempty.  Since $P_0(A)\in(0,1)$ and $A$ is
convex, $A$ has nonempty interior.  Taking a point of $\widehat A_\rho$ outside
that supporting halfspace and joining it to an interior point of $A$ gives
points of $\operatorname{int}(\widehat A_\rho)\setminus A$.  Thus
$\widehat A_\rho\setminus A$ contains a nonempty open set.  The density of
$N_2(0,\Omega_\rho)$ is everywhere positive, and therefore
\[
P_0(\widehat A_\rho)>P_0(A)~.
\]
\end{proof}

\begin{lemma}[Calibration of the competitor]\label{lem:competitor}
Suppose that $P_0(A)=1-\alpha$ and that $A$ has a finite-support direction
$a$ satisfying $a_2>a_1$.  Then there exists a finite number $b$ such that
\[
P_0\bigl(\widehat A_\rho
\cap\{x:T_{1,\rho}(x)\le b\}\bigr)=1-\alpha~.
\]
For any such $b$, define
\[
A':=\widehat A_\rho\cap\{x:T_{1,\rho}(x)\le b\}~.
\]
Then $A'$ is closed, convex, and coordinatewise lower, satisfies
$P_0(A')=1-\alpha$, and every nonzero finite-support direction of $A'$ lies in
$\mathcal F_\rho$.  Consequently, the test
$\phi_{A'}(x)=I\{x\notin A'\}$ has level $\alpha$, attained at $\mu=0$.
\end{lemma}

\begin{proof}
By Lemma~\ref{lem:hull},
\[
P_0(\widehat A_\rho)>P_0(A)=1-\alpha~.
\]
The map
\[
t\longmapsto
P_0\bigl(\widehat A_\rho
\cap\{x:T_{1,\rho}(x)\le t\}\bigr)
\]
is continuous because $T_{1,\rho}(X)$ has a nondegenerate normal distribution
under $P_0$.  It converges to zero as $t\to-\infty$ and to
$P_0(\widehat A_\rho)$ as $t\to\infty$.  Hence there exists a finite $b$ for
which the displayed probability equals $1-\alpha$.

The set $A'$ is closed and convex because it is the intersection of two sets
with these properties.  It is also coordinatewise lower.  Indeed,
\[
T_{1,\rho}(x)=\frac{x_1+x_2}{1+\rho}
\]
is coordinatewise increasing because $1+\rho>0$.  Its null probability equals
$1-\alpha$ by construction.

It remains to characterize the finite-support directions of $A'$.  We first
show that $A'$ is $\Pi_\rho$-lower.  Let $x\in A'$ and $p\in\Pi_\rho$.
Lemma~\ref{lem:hull} implies that $x-p\in\widehat A_\rho$.  Moreover,
$T_{1,\rho}(p)\ge0$ by the definition of $\Pi_\rho$, and therefore
\[
T_{1,\rho}(x-p)
=T_{1,\rho}(x)-T_{1,\rho}(p)
\le T_{1,\rho}(x)
\le b~.
\]
Thus $x-p\in A'$.

Now let $a\ne0$ satisfy $h_{A'}(a)<\infty$, and fix $x\in A'$.  Since $A'$ is
coordinatewise lower, $x-te_j\in A'$ for every $t\ge0$ and $j=1,2$.  If
$a_j<0$, then
\[
a'(x-te_j)=a'x-ta_j\longrightarrow+\infty~,
\]
contradicting $h_{A'}(a)<\infty$.  Hence $a_1\ge0$ and $a_2\ge0$.

Because $(1,-1)\in\Pi_\rho$ and $A'$ is $\Pi_\rho$-lower,
\[
x-t(1,-1)\in A'
\qquad\text{for every }t\ge0~.
\]
If $a_2>a_1$, then
\[
a'\{x-t(1,-1)\}
=a'x+t(a_2-a_1)
\longrightarrow+\infty~,
\]
again contradicting finite support.  Therefore $a_2\le a_1$.  We have shown
that every nonzero $a$ with $h_{A'}(a)<\infty$ satisfies
\[
0\le a_2\le a_1~,
\]
so $a\in\mathcal F_\rho$.

Finally, the complement of $A'$ is coordinatewise upper.  For every
$\mu\in\Lambda_0$, we have $\mu\le0$.  Coupling $X=\mu+Z$ with
$Z\sim N_2(0,\Omega_\rho)$ gives
\[
P_\mu(X\notin A')\le P_0(Z\notin A')=\alpha~.
\]
Equality holds at the finite null point $\mu=0$, which proves the last claim.
\end{proof}

\begin{lemma}[Transfer domination criterion]\label{lem:transfer}
Let $G$ and $D$ be Borel sets with
\[
P_0(G)=P_0(D)>0~.
\]
Suppose that, for every $g\in G$ and $d\in D$,
\[
T_{1,\rho}(g)>T_{1,\rho}(d)
\qquad\text{and}\qquad
T_{2,\rho}(g)\ge T_{2,\rho}(d)~.
\]
Then
\[
P_\mu(G)>P_\mu(D)
\qquad\text{for every }\mu\in\Lambda_1~.
\]
\end{lemma}

\begin{proof}
Let $m=P_0(G)=P_0(D)$.  Using the likelihood ratio in
\eqref{eq:LR-density},
\[
m\{P_\mu(G)-P_\mu(D)\}
=
\int_G\int_D
\{L_\mu(g)-L_\mu(d)\}\,dP_0(d)\,dP_0(g)~.
\]
For $\mu\in\Lambda_1$, the coefficient on $T_{1,\rho}$ in
\eqref{eq:LR-density} is $u=\mu_1>0$, while the coefficient on
$T_{2,\rho}$ is $v=\mu_1-\mu_2\ge0$.  The assumed ordering therefore gives
\[
L_\mu(g)>L_\mu(d)
\]
for every pair $(g,d)\in G\times D$.  Since $m>0$, the double integral is
strictly positive, which proves the result.
\end{proof}

\section*{3. Necessity}

It remains to prove that a bad finite normal generates sets $G$ and $D$
satisfying Lemma~\ref{lem:transfer}.  The proof is slicewise.  The first step is
to observe that, although the transfer cone $\Pi_\rho$ depends on $\rho$, its
lower hull does not.  Let
\[
K:=\operatorname{cone}\{(1,-1)\}~.
\]
Since $K\subseteq\Pi_\rho$, we have $A-K\subseteq A-\Pi_\rho$.  For the reverse
inclusion, recall that
\[
\Pi_\rho=\operatorname{cone}\{(1,\rho),(1,-1)\}~.
\]
If $\rho\ge0$, subtracting a nonnegative multiple of $(1,\rho)$ is a
coordinatewise downward move and is already absorbed by $A$.  If $\rho<0$,
then
\[
(1,\rho)=(-\rho)(1,-1)+(1+\rho)e_1~,
\]
where both coefficients are positive.  Subtracting the second term is again a
coordinatewise downward move.  Hence, in both cases,
\begin{equation}\label{eq:hull-reduction}
A-\Pi_\rho=A-K,
\qquad
\widehat A_\rho=\operatorname{cl}(A-K)~.
\end{equation}

For $s\in\mathbf R$, define the horizontal boundary of $A$ by
\[
\beta(s):=\sup\{x_1:(x_1,s)\in A\}\in[-\infty,+\infty]~.
\]
Since $A$ is closed, convex, and coordinatewise lower,
\[
A\cap\{x_2=s\}=\{(x_1,s):x_1\le\beta(s)\}~,
\]
with the usual interpretation when $\beta(s)=\pm\infty$.  The function
$\beta$ is concave in the extended-real sense.  Set
\[
m(s):=\beta(s)+s,
\qquad
M(s):=\sup_{r\le s}m(r)~.
\]
Thus $m(s)$ is the largest value of the unscaled diagonal score
$x_1+x_2=(1+\rho)T_{1,\rho}(x)$ on the horizontal slice of $A$, and $M(s)$ is
its running maximum from below.

\begin{lemma}[Slice formula for $\widehat A_\rho$]\label{lem:slice-formula}
For every $s\in\mathbf R$,
\[
\widehat A_\rho\cap\{x_2=s\}
=
\{(x_1,s):x_1\le M(s)-s\}~.
\]
Equivalently,
\[
(x_1,s)\in\widehat A_\rho
\quad\Longleftrightarrow\quad
x_1+s\le M(s)~.
\]
\end{lemma}

\begin{proof}
We first verify that the set
\[
C:=\{(x_1,s):x_1+s\le M(s)\}
\]
is closed.  The function $M$ is never $-\infty$.  To see this, fix
$(x_1^0,s_0)\in A$.  If $s\ge s_0$, then $M(s)\ge m(s_0)>-\infty$; if
$s<s_0$, coordinatewise lowerness gives $(x_1^0,s)\in A$, so
$M(s)\ge m(s)>-\infty$.

The function $M$ is nondecreasing.  It is also concave whenever it is finite.
Indeed, for $s_1,s_2\in\mathbf R$, choose $r_i\le s_i$ such that
$m(r_i)\ge M(s_i)-\varepsilon$, $i=1,2$.  For every $\lambda\in[0,1]$,
concavity of $m$ gives
\begin{align*}
M\{\lambda s_1+(1-\lambda)s_2\}
&\ge m\{\lambda r_1+(1-\lambda)r_2\}\\
&\ge \lambda M(s_1)+(1-\lambda)M(s_2)-\varepsilon~.
\end{align*}
Letting $\varepsilon\downarrow0$ proves the claim.

If $M$ takes the value $+\infty$ at some point, then it is identically
$+\infty$.  Monotonicity gives this conclusion to the right.  To see it to the
left, suppose $M(s_0)=+\infty$ and fix $t<s_0$.  Choose $r_n\le s_0$ with
$m(r_n)\to+\infty$.  If $r_n\le t$ along a subsequence, then $M(t)=+\infty$.
Otherwise, $t<r_n\le s_0$ eventually.  Since $A$ is coordinatewise lower,
$\beta(t)\ge\beta(r_n)$, and hence
\[
m(t)\ge m(r_n)+t-r_n
\ge m(r_n)+t-s_0\longrightarrow+\infty~.
\]
Thus $M(t)=+\infty$ in this case as well.  It follows that either
$M\equiv+\infty$, in which case $C=\mathbf R^2$, or $M$ is a finite concave
function on $\mathbf R$.  In the latter case $M$ is continuous.  Therefore $C$
is closed in either case.

Now suppose $(x_1,s)\in A-K$.  By \eqref{eq:hull-reduction}, there exist
$(y_1,r)\in A$ and $\tau\ge0$ such that
\[
(x_1,s)=(y_1,r)-\tau(1,-1)~.
\]
Hence $r=s-\tau\le s$ and
\[
x_1+s=y_1+r\le m(r)\le M(s)~.
\]
Thus $A-K\subseteq C$.  Since $C$ is closed,
$\widehat A_\rho=\operatorname{cl}(A-K)\subseteq C$.

Conversely, suppose $x_1+s<M(s)$.  Then there exists $r\le s$ such that
\[
m(r)>x_1+s~.
\]
Choose $y_1\le\beta(r)$ with $y_1+r>x_1+s$.  Then $(y_1,r)\in A$.  Since $A$
is coordinatewise lower, the point
\[
\widetilde y_1:=x_1+s-r<y_1
\]
also satisfies $(\widetilde y_1,r)\in A$.  Setting $\tau=s-r\ge0$ gives
\[
(x_1,s)=(\widetilde y_1,r)-\tau(1,-1)\in A-K~.
\]
Thus every point satisfying the strict inequality belongs to $A-K$.  If
$x_1+s=M(s)$, apply this argument to $(x_1-1/n,s)$ and let $n\to\infty$.
Since $\widehat A_\rho$ is closed, the point $(x_1,s)$ also belongs to
$\widehat A_\rho$.  This proves the formula.
\end{proof}

For $\rho<0$, write
\[
q_\rho:=(1,\rho)~,
\]
and say that a set $C$ is $q_\rho$-lower if
\[
x\in C,\quad t\ge0
\quad\Longrightarrow\quad
x-tq_\rho\in C~.
\]
This property is automatic when $\rho\ge0$ and $C$ is coordinatewise lower,
but not when $\rho<0$: subtracting $q_\rho$ lowers the first coordinate and
raises the second.  The next two lemmas handle the two possibilities.

\begin{lemma}[Conditional rearrangement]\label{lem:rearrangement}
Suppose $\rho<0$.  If $A$ is not $q_\rho$-lower, then there exists a Borel
acceptance region $A^{\mathrm r}$ such that the test
$\phi_{A^{\mathrm r}}=I\{X\notin A^{\mathrm r}\}$ has level $\alpha$ and
\[
P_\mu(A^{\mathrm r})<P_\mu(A)
\qquad\text{for every }\mu\in\Lambda_1~.
\]
Consequently, $\phi_{A^{\mathrm r}}$ strictly dominates $\phi_A$.
\end{lemma}

\begin{proof}
Introduce the coordinates
\[
Z=X_1,
\qquad
E=X_2-\rho X_1~.
\]
Writing $\mu=(u,u-v)$, where $v\ge0$ on $\Lambda$, we have
\[
Z\sim N(u,1),
\qquad
E\sim N((1-\rho)u-v,1-\rho^2),
\qquad
Z\perp E~.
\]
For $e\in\mathbf R$, let
\[
A_e:=\{z:(z,e+\rho z)\in A\},
\qquad
p(e):=\Phi(A_e),
\]
where $\Phi$ denotes standard normal probability.  By Fubini's theorem,
$p$ is measurable.  Moreover, if $e_1\le e_2$, coordinatewise lowerness of
$A$ gives $A_{e_2}\subseteq A_{e_1}$.  Hence $p$ is nonincreasing.

Let
\[
c(e):=\Phi^{-1}\{p(e)\},
\]
with the conventions $c(e)=-\infty$ when $p(e)=0$ and $c(e)=+\infty$ when
$p(e)=1$, and define
\begin{equation}\label{eq:rearranged-set}
A^{\mathrm r}
:=\{x:x_1\le c(x_2-\rho x_1)\}~.
\end{equation}
The function $c$ is measurable and nonincreasing.  Thus, in $(z,e)$
coordinates, $A^{\mathrm r}=\{(z,e):z\le c(e)\}$ is Borel and
coordinatewise lower.

Under $P_0$, $Z$ and $E$ are independent and the conditional standard normal
mass of every section is preserved.  It follows that
\[
P_0(A^{\mathrm r})=P_0(A)=1-\alpha~.
\]
If $\mu\in\Lambda_0$, then $u\le0$, $v\ge0$, and
\[
\bigl(u,(1-\rho)u-v\bigr)\le(0,0)~.
\]
Since $A^{\mathrm r}$ is lower in $(z,e)$ coordinates, translation coupling
therefore gives
\[
P_\mu(A^{\mathrm r})\ge P_0(A^{\mathrm r})=1-\alpha~.
\]
The rearranged test consequently has level $\alpha$, with equality at
$\mu=0$.

Now fix $\mu\in\Lambda_1$, so that $u>0$.  Conditional on $E=e$, the
likelihood ratio of $N(u,1)$ relative to $N(0,1)$ is strictly increasing in
$z$.  Among all subsets of the real line having standard normal probability
$p(e)$, the left halfline $(-\infty,c(e)]$ therefore minimizes their
$N(u,1)$ probability.  Applying this conditional Neyman--Pearson inequality
and integrating over $E$ yields
\begin{equation}\label{eq:weak-rearrangement}
P_\mu(A^{\mathrm r})\le P_\mu(A)~.
\end{equation}

It remains to prove strictness.  Because $A$ is not $q_\rho$-lower, there are
$x\in A$ and $t>0$ for which $x-tq_\rho\notin A$.  The set $A$ has nonempty
interior.  By moving $x$ slightly toward an interior point of $A$, if
necessary, we may find $x^0\in\operatorname{int}(A)$ and $t_0>0$ such that
$x^0-t_0q_\rho\notin A$.  Closedness of $A$ then gives small open
neighborhoods of these two points, the first contained in $A$ and the second
contained in its complement.  Under the invertible map
\[
H(x)=(x_1,x_2-\rho x_1),
\]
the two centers have the same second coordinate and their first coordinates
differ by $t_0$.  The images of the neighborhoods consequently contain two
rectangles with a common nonempty interval of $e$ values.  For every $e$ in a
smaller such interval, $A_e$ contains an interval lying strictly to the right
of an interval it omits.  Thus $A_e$ is not a left halfline modulo standard
normal measure.  The conditional Neyman--Pearson inequality is strict for all
these $e$.  Since the density of $E$ is everywhere positive, integration
makes \eqref{eq:weak-rearrangement} strict for every $u>0$, and hence for every
$\mu\in\Lambda_1$.
\end{proof}

\begin{lemma}[Gain-loss separation]\label{lem:gain-loss}
Let $A'$ be defined as in Lemma~\ref{lem:competitor}, and define the gain and
loss sets
\[
G:=A\setminus A'
=A\cap\{T_{1,\rho}>b\},
\qquad
D:=A'\setminus A~.
\]
Suppose either that $0\le\rho<1$, or that $-1<\rho<0$ and $A$ is
$q_\rho$-lower.  Then
\[
P_0(G)=P_0(D)>0~.
\]
Moreover, for every $g\in G$ and every $d\in D$,
\[
T_{1,\rho}(g)>T_{1,\rho}(d)
\qquad\text{and}\qquad
T_{2,\rho}(g)\ge T_{2,\rho}(d)~.
\]
\end{lemma}

\begin{proof}
Since $A\subseteq\widehat A_\rho$ and
$A'=\widehat A_\rho\cap\{T_{1,\rho}\le b\}$,
\[
A\setminus A'=A\cap\{T_{1,\rho}>b\}=G~.
\]
Also $D=A'\setminus A$.  Because $P_0(A)=P_0(A')=1-\alpha$,
\[
P_0(G)=P_0(D)~.
\]
This common mass is positive.  Otherwise $A$ and $A'$ would be equal up to a
Gaussian null set.  Since both sets are closed, convex, and have positive
Gaussian probability, they have nonempty interiors.  If the sets were not
equal, joining a point in one set but not the other to an interior point of the
first set would produce a nonempty open subset of their symmetric difference.
The Gaussian density is everywhere positive, so the symmetric difference could
not have probability zero.  Hence null symmetric difference would imply
$A=A'$.  This is impossible because Lemma~\ref{lem:competitor} shows that every
finite-support direction of $A'$ lies in $\mathcal F_\rho$, whereas $A$ has a
finite-support direction outside $\mathcal F_\rho$.

We now prove the ordering.  Write
\[
B_\rho:=(1+\rho)b~,
\]
so that $T_{1,\rho}(x)\le b$ is equivalent to
$x_1+x_2\le B_\rho$.  On the slice $x_2=s$, the gain set is
\[
G_s
:=\{x_1:(x_1,s)\in G\}
=
(B_\rho-s,\beta(s)]~.
\]
Thus
\begin{equation}\label{eq:gain-slice}
G_s\ne\varnothing
\quad\Longleftrightarrow\quad
m(s)>B_\rho~.
\end{equation}
Using Lemma~\ref{lem:slice-formula}, the loss slice is
\[
D_s
:=\{x_1:(x_1,s)\in D\}
=
(\beta(s),\min\{M(s),B_\rho\}-s]~.
\]
Consequently,
\begin{equation}\label{eq:loss-slice}
D_s\ne\varnothing
\quad\Longleftrightarrow\quad
M(s)>m(s)
\quad\text{and}\quad
m(s)<B_\rho~.
\end{equation}

Let
\[
s_+:=\sup\{s:m(s)>B_\rho\}~.
\]
The set in braces is nonempty because $P_0(G)>0$.  We claim that
\[
G_s\ne\varnothing\quad\Longrightarrow\quad s\le s_+,
\]
and
\[
D_s\ne\varnothing\quad\Longrightarrow\quad s\ge s_+~.
\]
The first implication follows from the definition of $s_+$ and
\eqref{eq:gain-slice}.  For the second, suppose toward a contradiction that
$D_s\ne\varnothing$ for some $s<s_+$.  By \eqref{eq:loss-slice},
$M(s)>m(s)$, so there is $q\le s$ such that
\[
m(q)>m(s)~.
\]
Necessarily $q<s$.  Since $s<s_+$, there is also $r>s$ such that
\[
m(r)>B_\rho>m(s)~.
\]
Because $q<s<r$ and $m$ is concave, for the unique $\lambda\in(0,1)$ satisfying
$s=\lambda q+(1-\lambda)r$ we obtain
\[
m(s)
\ge
\lambda m(q)+(1-\lambda)m(r)
>
m(s)~,
\]
a contradiction.  Hence no loss slice lies below $s_+$.  Since $D$ is nonempty,
this also shows that $s_+$ is finite.

It follows that every $g\in G$ and $d\in D$ satisfy
\[
g_2\le s_+\le d_2~.
\]
In fact, $g_2<d_2$, because a gain slice requires $m(s)>B_\rho$, whereas a
loss slice requires $m(s)<B_\rho$.
By definition of the gain and loss sets,
\[
T_{1,\rho}(g)>b\ge T_{1,\rho}(d)~.
\]
If $\rho\ge0$, the identity
\[
T_{2,\rho}(x)
=\frac{\rho T_{1,\rho}(x)-x_2}{1-\rho}
\]
gives
\[
T_{2,\rho}(g)-T_{2,\rho}(d)
=
\frac{
\rho\{T_{1,\rho}(g)-T_{1,\rho}(d)\}+d_2-g_2
}{1-\rho}
\ge0~,
\]
which proves the result in this case.

Suppose instead that $\rho<0$ and $A$ is $q_\rho$-lower.  If $s<t$ and
$(x_1,s)\in A$, set $\tau=(t-s)/(-\rho)$.  Then
\[
(x_1,s)-\tau q_\rho=(x_1-\tau,t)\in A~.
\]
It follows that
\[
\beta(t)\ge\beta(s)-\frac{t-s}{-\rho}
\]
whenever the relevant boundaries are finite, and hence
\begin{equation}\label{eq:slope-bound}
m(s)-m(t)
\le
\frac{1+\rho}{-\rho}(t-s)~.
\end{equation}
For $g\in G$ and $d\in D$, put $s=g_2$ and $t=d_2$.  We have $s<t$,
and the boundaries at $s$ and $t$ are finite.  Indeed, a gain point makes the
slice at $s$ nonempty, while $q_\rho$-lowerness carries such a point to height
$t$; moreover, $\beta(s)=+\infty$ would force $\beta(t)=+\infty$, contrary to
the existence of a loss point at $t$.  Therefore
\[
g_1+s\le m(s),
\qquad
d_1+t>m(t),
\]
and therefore, by \eqref{eq:slope-bound},
\[
(g_1+g_2)-(d_1+d_2)
<\frac{1+\rho}{-\rho}(d_2-g_2)~.
\]
Since $\rho<0$, multiplication reverses this strict inequality.  Consequently,
\begin{align*}
(1-\rho^2)\{T_{2,\rho}(g)-T_{2,\rho}(d)\}
&=\rho\{(g_1+g_2)-(d_1+d_2)\}
 +(1+\rho)(d_2-g_2)\\
&>0~.
\end{align*}
This proves the asserted pairwise ordering in both cases.
\end{proof}

\begin{lemma}[Necessity of feasible normals]\label{lem:necessity}
If $A$ has a finite-support direction $a$ with $a_2>a_1$, then $\phi_A$ is
inadmissible.  More precisely, there exists a level-$\alpha$ test $\phi^*$
such that
\[
P_\mu(\phi^*=1)>P_\mu(\phi_A=1)
\qquad\text{for every }\mu\in\Lambda_1~.
\]
\end{lemma}

\begin{proof}
First suppose $\rho\ge0$, and construct $A'$ as in
Lemma~\ref{lem:competitor}.  That lemma shows that $\phi_{A'}$ has level
$\alpha$.  Writing $G=A\setminus A'$ and $D=A'\setminus A$, the difference in
power is
\[
P_\mu(X\notin A')-P_\mu(X\notin A)
=P_\mu(G)-P_\mu(D)~.
\]
Lemma~\ref{lem:gain-loss} verifies the hypotheses of the transfer criterion,
Lemma~\ref{lem:transfer}.  Hence this difference is strictly positive for
every $\mu\in\Lambda_1$.

Now suppose $\rho<0$.  If $A$ is not $q_\rho$-lower,
Lemma~\ref{lem:rearrangement} supplies a level-$\alpha$ test that strictly
dominates $\phi_A$ at every $\mu\in\Lambda_1$.  If $A$ is $q_\rho$-lower,
construct $A'$ as in Lemma~\ref{lem:competitor}.  The negative-correlation
part of Lemma~\ref{lem:gain-loss} again verifies the hypotheses of
Lemma~\ref{lem:transfer}, so $\phi_{A'}$ strictly dominates $\phi_A$ on
$\Lambda_1$.  These cases exhaust $-1<\rho<1$.
\end{proof}

\begin{proof}[Proof of Theorem~\ref{thm:iff}]
If all finite normals of $A$ lie in $\mathcal F_\rho$,
Lemma~\ref{lem:sufficiency} gives $\alpha$-admissibility against all
level-$\alpha$ tests.

Conversely, suppose the finite-normal condition fails.  Since $A$ is
coordinatewise lower, every finite-support direction is nonnegative.  Thus
there exists $a\ne0$ such that
\[
h_A(a)<\infty,
\qquad
0\le a_1<a_2~.
\]
Lemma~\ref{lem:necessity} constructs a level-$\alpha$ test that strictly
dominates $\phi_A$ at every $\mu\in\Lambda_1$.  Therefore $\phi_A$ is
inadmissible.
\end{proof}

\begin{remark}[Why the identity case needs no extra step]
When $\rho=0$, the score $T_{2,0}(x)=-x_2$.  The vertical separation
$g_2<d_2$ therefore gives the required score ordering directly.  Equivalently,
$q_0=(1,0)$, and coordinatewise lowerness already implies $q_0$-lowerness.
For $\rho<0$, neither fact remains true: $T_{2,\rho}$ contains a negative
multiple of $T_{1,\rho}$, and subtracting $q_\rho$ raises the second
coordinate.  The rearrangement lemma handles acceptance regions that fail
this new directional monotonicity; the slope bound in
Lemma~\ref{lem:gain-loss} handles those that satisfy it.  This is precisely
the additional work required beyond the identity-covariance proof.
\end{remark}

\end{document}
